% !TEX program = pdflatex
% Build:
%   cd docs/notes && pdflatex scikit-cuda-dependency.tex
% or:
%   latexmk -pdf scikit-cuda-dependency.tex
\documentclass[11pt,a4paper]{article}

\usepackage[margin=1in]{geometry}
\usepackage[T1]{fontenc}
\usepackage{lmodern}
\usepackage{microtype}
\usepackage{xcolor}
\usepackage{hyperref}
\usepackage{listings}
\usepackage{enumitem}
\usepackage{titlesec}

\hypersetup{
    colorlinks=true,
    linkcolor=black,
    urlcolor=blue!60!black,
    citecolor=black,
}

\definecolor{codebg}{RGB}{245,245,245}
\definecolor{codecomment}{RGB}{106,153,85}
\definecolor{codekw}{RGB}{0,0,180}
\definecolor{codestr}{RGB}{163,21,21}

\lstset{
    backgroundcolor=\color{codebg},
    basicstyle=\ttfamily\small,
    keywordstyle=\color{codekw}\bfseries,
    commentstyle=\color{codecomment}\itshape,
    stringstyle=\color{codestr},
    breaklines=true,
    breakatwhitespace=true,
    showstringspaces=false,
    frame=single,
    framesep=4pt,
    rulecolor=\color{codebg},
    xleftmargin=8pt,
    xrightmargin=8pt,
    aboveskip=8pt,
    belowskip=8pt,
}

\lstdefinelanguage{pyplus}{
    language=Python,
    morekeywords={cufft,skcuda,gpuarray,np},
}

\titleformat{\section}{\large\bfseries}{\thesection}{0.6em}{}
\titleformat{\subsection}{\normalsize\bfseries}{\thesubsection}{0.6em}{}

\title{Technical Note: \texttt{scikit-cuda} as a Dependency of \texttt{cuvarbase}}
\author{cuvarbase fork (\texttt{WilliamCFong/cuvarbase})}
\date{\today}

\begin{document}
\maketitle

\begin{abstract}
\noindent
\texttt{cuvarbase} is a PyCUDA-backed library of GPU period-finding tools for
astronomical time series. It depends on \texttt{scikit-cuda} for one thing
only: GPU FFT plans and the inverse-FFT call inside the adjoint NFFT used by
the Lomb--Scargle periodogram. \texttt{scikit-cuda} has not had a release since
June 2019 and is, in practice, abandoned. This note documents (a)~exactly how
\texttt{cuvarbase} consumes \texttt{scikit-cuda}, (b)~the concrete failures
caused by upstream's lack of maintenance against modern NumPy and Setuptools,
(c)~the workaround currently applied in this fork, and (d)~a recommended
migration path.
\end{abstract}

\section{Background}

\subsection{What \texttt{cuvarbase} is}
\texttt{cuvarbase} implements several time-series period finders on the GPU:
generalized Lomb--Scargle (GLS), Box-Least-Squares (BLS), Conditional Entropy
(CE), Phase-Dispersion Minimization (PDM2), and an adjoint Non-equispaced FFT
(NFFT) that the GLS path consumes. Kernels live in
\texttt{src/cuvarbase/kernels/*.cu} and are compiled at runtime via PyCUDA's
\texttt{SourceModule}. PyCUDA itself supplies device memory, streams, and the
JIT compiler --- but \emph{not} an FFT.

\subsection{What \texttt{scikit-cuda} is}
\texttt{scikit-cuda} (\href{https://github.com/lebedov/scikit-cuda}{\texttt{lebedov/scikit-cuda}})
is a Python wrapper around NVIDIA's CUDA math libraries (cuBLAS, cuFFT,
cuSOLVER, cuRAND, etc.). For the FFT in particular, \texttt{skcuda.fft} wraps
cuFFT plans and execution functions so callers can stay in Python and exchange
PyCUDA \texttt{GPUArray}s.

\subsection{Why \texttt{cuvarbase} pulls it in}
The adjoint NFFT in \texttt{src/cuvarbase/cunfft.py} needs an inverse FFT on a
gridded, oversampled spectrum after the kernel-spreading step. Writing a cuFFT
binding by hand was avoided in 2017 by reusing
\texttt{skcuda.fft.Plan}/\texttt{ifft}.

\section{How \texttt{cuvarbase} actually uses \texttt{scikit-cuda}}

The surface area is small. Across the entire package and test suite there are
exactly two call sites for two APIs:

\begin{lstlisting}[language=pyplus,caption={\texttt{src/cuvarbase/cunfft.py}}]
import skcuda.fft as cufft
# ...
self.cu_plan = cufft.Plan(self.n, self.complex_type, self.complex_type,
                          stream=self.stream)
# ...
cufft.ifft(memory.ghat_g, memory.ghat_g, memory.cu_plan)
\end{lstlisting}

\begin{lstlisting}[language=pyplus,caption={\texttt{tests/test\_nfft.py}}]
import skcuda.fft as cufft
# ...
plan = cufft.Plan(len(y), np.complex128, np.complex128)
cufft.ifft(yg, yghat, plan)
\end{lstlisting}

That is the entire dependency: a constructor (\texttt{Plan}) and an in-place
inverse-FFT call (\texttt{ifft}). No cuBLAS, cuSOLVER, cuRAND, or other
\texttt{skcuda.*} subpackage is touched. Everything else is unaffected.

\section{Maintenance status of \texttt{scikit-cuda}}

\begin{itemize}[leftmargin=1.4em]
    \item \textbf{Last PyPI release:} 0.5.3 (May 2019). No release in the
    intervening $\sim$6 years (as of \today).
    \item \textbf{Upstream activity:} the GitHub repository
    \texttt{lebedov/scikit-cuda} accumulates issues and PRs without merges. It
    is, by any reasonable definition, archived in spirit if not in flag.
    \item \textbf{Wheel availability:} two \texttt{py2.py3-none-any} wheels
    exist on PyPI for 0.5.2/0.5.3 plus three sdists. The pure-Python wheels
    install cleanly; the runtime requirements (CUDA, PyCUDA, recent NumPy) do
    not.
\end{itemize}

\section{Concrete breakage caused by lack of maintenance}

\subsection{NumPy 1.25 removed \texttt{np.typeDict}}
\label{sec:typedict}
\texttt{skcuda/misc.py:637} contains:

\begin{lstlisting}[language=pyplus]
num_types = [np.typeDict[t] for t in
             ('float32', 'float64', 'complex64', 'complex128')]
\end{lstlisting}

\noindent
\texttt{numpy.typeDict} was deprecated in NumPy~1.20 and removed in
NumPy~1.25. Any \texttt{import skcuda.fft} (or any \texttt{from skcuda import
*}) under NumPy~$\geq 1.25$ raises:

\begin{lstlisting}[basicstyle=\ttfamily\small]
AttributeError: module 'numpy' has no attribute 'typeDict'
\end{lstlisting}

\noindent
On a fresh \texttt{nox -s tests} install with no NumPy pin, pip resolves to
NumPy 2.x and both \texttt{tests/test\_lombscargle.py} and
\texttt{tests/test\_nfft.py} fail at \emph{collection time} --- the test files
cannot even be imported.

\subsection{\texttt{pkg\_resources} deprecation}
\texttt{skcuda/\_\_init\_\_.py} declares its namespace via
\texttt{pkg\_resources.declare\_namespace}. Setuptools is in the process of
removing \texttt{pkg\_resources} (slated for as early as 2025-11-30 per the
deprecation warning emitted on every import). Once the API is gone,
\texttt{import skcuda} fails outright. The project is currently emitting:

\begin{lstlisting}[basicstyle=\ttfamily\small]
UserWarning: pkg_resources is deprecated as an API.
DeprecationWarning: Deprecated call to
    pkg_resources.declare_namespace('skcuda').
\end{lstlisting}

\subsection{PyCUDA version drift}
\texttt{cuvarbase}'s pin already excludes \texttt{pycuda~==~2024.1.2} (a
known-bad release), but scikit-cuda has not been validated against any PyCUDA
released in the last several years. Future PyCUDA API changes carry no
guarantee of an upstream skcuda fix.

\section{Current workaround in this fork}

PR \#2 caps NumPy below the breaking version in \texttt{pyproject.toml}:

\begin{lstlisting}[language=pyplus,caption={\texttt{pyproject.toml}}]
dependencies = [
    # scikit-cuda 0.5.3 still uses `np.typeDict` (skcuda/misc.py:637),
    # which was removed in NumPy 1.25. Cap numpy until skcuda releases
    # a fix.
    "numpy>=1.6,<1.25",
    "scipy",
    "pycuda>=2017.1.1,!=2024.1.2",
    "scikit-cuda",
]
\end{lstlisting}

\noindent
This restores test collection. It is a stop-gap with two costs:

\begin{enumerate}[leftmargin=1.6em]
    \item NumPy 1.24 is the last version under the cap. Subsequent NumPy
    improvements (performance, new APIs, native macOS arm64 support, NumPy 2
    typing) are inaccessible to consumers of \texttt{cuvarbase}.
    \item NumPy 1.24 has no Python 3.12 wheels, so the practical Python
    upper bound is 3.11 even though \texttt{requires-python} still says
    \texttt{>=3.9} and the classifiers list 3.12.
\end{enumerate}

The pin does \emph{nothing} for the \texttt{pkg\_resources} timebomb in
section~\textsection 4.2.

\section{Recommended migration}

The skcuda surface area in \texttt{cuvarbase} is small enough --- one
\texttt{Plan} object, one \texttt{ifft} call --- that direct replacement is
straightforward. Three viable replacements, in rough order of fit:

\begin{description}[leftmargin=1.6em,style=nextline]
    \item[\texttt{cuda-python} (\texttt{nvidia-cuda-python})]
    NVIDIA's official low-level Python bindings. Includes cuFFT bindings via
    \texttt{cuda.bindings.cufft}. Supported, versioned with the CUDA toolkit,
    and the canonical answer going forward. Requires writing a thin wrapper
    (plan-handle creation, exec call, destroy) but the control flow is a
    direct map from the existing two call sites.

    \item[\texttt{nvmath-python}]
    NVIDIA's high-level math-libraries package (cuFFT, cuBLAS, cuSPARSE,
    cuSOLVER) layered on top of \texttt{cuda-python}. Closer ergonomically to
    \texttt{skcuda.fft}, but newer and dependency-heavy.

    \item[\texttt{cupy}]
    \texttt{cupy.fft} wraps cuFFT and integrates with PyCUDA arrays via the
    \texttt{\_\_cuda\_array\_interface\_\_} protocol. Pulls in CuPy as a
    dependency, which is a sizeable addition for what amounts to two function
    calls.
\end{description}

\subsection*{Estimated effort}
\begin{itemize}[leftmargin=1.6em]
    \item \texttt{src/cuvarbase/cunfft.py}: replace the \texttt{Plan}
    construction in \texttt{NFFTMemory.allocate\_grid} and the
    \texttt{cufft.ifft} call in \texttt{nfft\_adjoint\_async}.
    \item \texttt{tests/test\_nfft.py}: replace the equivalent two lines used
    in the CPU-vs-GPU comparison fixture.
    \item Update \texttt{pyproject.toml}: drop \texttt{scikit-cuda}, drop the
    \texttt{numpy<1.25} cap, add the chosen replacement.
    \item Update CLAUDE.md and \texttt{INSTALL.rst}: scikit-cuda is no longer
    a runtime dependency.
\end{itemize}

\noindent
A reasonable bound: a single afternoon of focused work plus a re-run of the
GPU test suite. The migration is a strict removal of risk; nothing is lost.

\section{Action items}

\begin{enumerate}[leftmargin=1.6em]
    \item \textbf{Short term:} keep the \texttt{numpy<1.25} cap from PR \#2 in
    place. Document it (this note).
    \item \textbf{Medium term:} prototype the \texttt{cuda-python} cuFFT path
    in a feature branch, validated against \texttt{tests/test\_nfft.py}.
    \item \textbf{Long term:} delete \texttt{scikit-cuda} from
    \texttt{dependencies}, lift the NumPy cap, restore Python 3.12 to the
    supported matrix, and remove this note.
\end{enumerate}

\end{document}
